\documentclass[12pt,a4paper]{ctexart}

\usepackage[a4paper,margin=1.8cm]{geometry}
\usepackage{amsmath,amssymb}
\usepackage{enumitem}
\usepackage{xcolor}
\usepackage{listings}
\usepackage{tcolorbox}
\usepackage{titlesec}
\usepackage{setspace}
\usepackage{hyperref}
% 水印部分
% 直接挂载到 LaTeX 的前景(foreground)渲染层
\AddToHook{shipout/foreground}{
    \begin{tikzpicture}[remember picture, overlay]
        % 在页面正中心放置置顶水印
        \node[text=blue, opacity=0.2, scale=5, rotate=30] 
        at (current page.center) {fifine专升本};
    \end{tikzpicture}
}
\setstretch{1.2}
\setlist[enumerate]{itemsep=0.6em, topsep=0.4em}

\titleformat{\section}{\Large\bfseries}{}{0em}{}
\titleformat{\subsection}{\large\bfseries}{}{0em}{}

\lstset{
    language=C,
    basicstyle=\ttfamily\small,
    keywordstyle=\color{blue},
    commentstyle=\color{gray},
    stringstyle=\color{red!70!black},
    numbers=left,
    numberstyle=\tiny\color{gray},
    stepnumber=1,
    numbersep=8pt,
    frame=single,
    breaklines=true,
    columns=flexible,
    showstringspaces=false,
    tabsize=4
}

\newtcolorbox{coverbox}{
    colback=gray!5,
    colframe=black!60,
    boxrule=0.8pt,
    arc=2mm,
    left=2mm,right=2mm,top=2mm,bottom=2mm
}

\newtcolorbox{tipbox}{
    colback=white,
    colframe=black!40,
    boxrule=0.6pt,
    arc=1.5mm,
    left=2mm,right=2mm,top=1.5mm,bottom=1.5mm
}

\begin{document}

\begin{center}
    {\LARGE \textbf{专升本 C 语言小红书发题题库（第一阶段）}}\\[0.6em]
    {\large 适用方向：专升本学生｜题目型内容｜仅发题版｜可用于早题晚解析}\\[0.8em]
    {\normalsize 使用方式：每一节可单独作为 1 篇小红书笔记发布}
\end{center}

\vspace{1em}

\section*{使用说明}
本题库为“小红书发题版”，每组题控制在 3～4 题，便于做成图文轮播或单篇发题笔记。\\
本文件 \textbf{不含答案、不含解析}，便于你后续单独制作“晚间解析版”。

\newpage

%========================
\section{笔记 1：数据类型 + 表达式入门母题}
\begin{coverbox}
\textbf{封面标题建议：}\\
专升本C语言｜数据类型和表达式，3题测出你基础稳不稳\\
\textbf{副标题建议：}\\
第1题会了，后面两题就顺了
\end{coverbox}

\subsection*{题目 1}
下列关于 C 语言用户标识符的叙述中，正确的是
\begin{enumerate}[label=\Alph*.]
    \item 用户标识符中可以出现下划线和减号
    \item 用户标识符中不可以出现减号，但可以出现下划线
    \item 用户标识符中可以出现下划线，但不可以放在开头
    \item 用户标识符中可以出现数字，它们也都可以放在开头
\end{enumerate}

\subsection*{题目 2}
请选出可用作 C 语言用户标识符的是
\begin{enumerate}[label=\Alph*.]
    \item \texttt{void, define, WORD}
    \item \texttt{a3\_b3, \_123, IF}
    \item \texttt{FOR, --abc, Case}
    \item \texttt{2a, Do, Sizeof}
\end{enumerate}

\subsection*{题目 3}
以下叙述中正确的是
\begin{enumerate}[label=\Alph*.]
    \item 构成 C 程序的基本单位是函数
    \item 可以在一个函数中定义另一个函数
    \item \texttt{main()} 函数必须放在其他函数之前
    \item C 函数定义的格式固定为 K\&R 格式
\end{enumerate}

\subsection*{题目 4}
一个 C 语言程序是由
\begin{enumerate}[label=\Alph*.]
    \item 一个主程序和若干子程序组成
    \item 函数组成
    \item 若干过程组成
    \item 若干子程序组成
\end{enumerate}

\newpage

%========================
\section{笔记 2：表达式最容易丢分的地方}
\begin{coverbox}
\textbf{封面标题建议：}\\
专升本C语言｜表达式这3题，很多人第一眼就做错\\
\textbf{副标题建议：}\\
别死算，先抓类型和运算规则
\end{coverbox}

\subsection*{题目 1}
在 16 位 C 编译系统上，若定义 \texttt{long a;}，则能给 \texttt{a} 赋值 40000 的正确语句是
\begin{enumerate}[label=\Alph*.]
    \item \texttt{a=20000+20000;}
    \item \texttt{a=4000*10;}
    \item \texttt{a=30000+10000;}
    \item \texttt{a=4000L*10L;}
\end{enumerate}

\subsection*{题目 2}
若 \texttt{a} 是 \texttt{int} 型变量，且 \texttt{a=5}，则表达式
\[
(a+100)\%2 + a/2
\]
的值为
\begin{enumerate}[label=\Alph*.]
    \item 52
    \item 51
    \item 53
    \item 50
\end{enumerate}

\subsection*{题目 3}
设有定义：\texttt{int a=12;}，则执行完语句
\[
a += a -= a*a;
\]
后，\texttt{a} 的值是
\begin{enumerate}[label=\Alph*.]
    \item 552
    \item 264
    \item 144
    \item -264
\end{enumerate}

\newpage

%========================
\section{笔记 3：if-else 嵌套判断母题}
\begin{coverbox}
\textbf{封面标题建议：}\\
专升本C语言｜if-else 嵌套判断，3题吃透\\
\textbf{副标题建议：}\\
这类题不是难，是很多同学配对关系总看错
\end{coverbox}

\subsection*{题目 1}
以下程序的输出结果是
\begin{lstlisting}
#include <stdio.h>
void main()
{
    int a=1,b=4,c=2;
    if(a<b)
        if(b<0) c=0;
        else c++;
    printf("%d\n",c);
}
\end{lstlisting}

\begin{enumerate}[label=\Alph*.]
    \item 0
    \item 1
    \item 2
    \item 3
\end{enumerate}

\subsection*{题目 2}
有以下程序
\begin{lstlisting}
#include <stdio.h>
void main()
{
    int a=1,b=0;
    if(!a) b++;
    else if(a==0)
        if(a) b++;
        else b++;
    printf("%d\n",b);
}
\end{lstlisting}

程序运行后的输出结果是
\begin{enumerate}[label=\Alph*.]
    \item 0
    \item 1
    \item 2
    \item 3
\end{enumerate}

\subsection*{题目 3}
以下程序的输出结果是
\begin{lstlisting}
#include <stdio.h>
void main()
{
    int a=0,b=1,c=0,d=20;
    if(a) d=d-10;
    else if(!b)
        if(!c) d=15;
        else d=25;
    printf("%d\n",d);
}
\end{lstlisting}

\begin{enumerate}[label=\Alph*.]
    \item 10
    \item 20
    \item 30
    \item 程序进入死循环
\end{enumerate}

\newpage

%========================
\section{笔记 4：switch 和 break 高频易错题}
\begin{coverbox}
\textbf{封面标题建议：}\\
专升本C语言｜switch 不会丢分？先做这3题\\
\textbf{副标题建议：}\\
真正难点不是语法，是执行流程
\end{coverbox}

\subsection*{题目 1}
下面程序的运行结果是
\begin{lstlisting}
#include <stdio.h>
void main()
{
    int x=1,y=0,a=0,b=0;
    switch(x)
    {
    case 1:
        switch(y)
        {
        case 0: a++; break;
        case 1: b++; break;
        }
    case 2: a++; b++; break;
    case 3: a++; b++;
    }
    printf("a=%d,b=%d\n",a,b);
}
\end{lstlisting}

\begin{enumerate}[label=\Alph*.]
    \item \texttt{a=2,b=1}
    \item \texttt{a=1,b=1}
    \item \texttt{a=1,b=0}
    \item \texttt{a=2,b=2}
\end{enumerate}

\subsection*{题目 2}
\texttt{continue} 语句的作用是
\begin{enumerate}[label=\Alph*.]
    \item 结束整个循环
    \item 结束本次循环，进入下一次循环
    \item 跳出 \texttt{switch} 语句
    \item 结束程序
\end{enumerate}

\subsection*{题目 3}
\texttt{break} 语句通常用在\_\_\_\_\_\_语句和循环语句中。
\begin{enumerate}[label=\Alph*.]
    \item \texttt{if}
    \item \texttt{switch}
    \item \texttt{for}
    \item \texttt{while}
\end{enumerate}

\newpage

%========================
\section{笔记 5：数组和指针的第一组母题}
\begin{coverbox}
\textbf{封面标题建议：}\\
专升本C语言｜数组和指针到底会不会？先做这4题\\
\textbf{副标题建议：}\\
这组题最适合区分“会背”和“真会做”
\end{coverbox}

\subsection*{题目 1}
若有定义：
\[
\texttt{int a[10], *p=a;}
\]
则 \texttt{*(p+5)} 表示
\begin{enumerate}[label=\Alph*.]
    \item \texttt{a[5]} 的地址
    \item \texttt{a[5]} 的值
    \item \texttt{a[6]} 的值
    \item \texttt{a} 的地址
\end{enumerate}

\subsection*{题目 2}
以下程序的输出结果是
\begin{lstlisting}
#include <stdio.h>
void main()
{
    int a[10]={1,2,3,4,5,6,7,8,9,10},*p=&a[3],b;
    b=*(p+2);
    printf("%d\n",b);
}
\end{lstlisting}

\begin{enumerate}[label=\Alph*.]
    \item 6
    \item 7
    \item 8
    \item 9
\end{enumerate}

\subsection*{题目 3}
若有以下定义和语句：
\[
	\texttt{int a[10]=\{1,2,3,4,5,6,7,8,9,10\}, *p=a;}
\]
则值为 6 的表达式是
\begin{enumerate}[label=\Alph*.]
	\item \texttt{*p+6}
	\item \texttt{*p++}
	\item \texttt{p+6}
	\item \texttt{p+=5,*p}
\end{enumerate}

\subsection*{题目 4}
若已有定义：\texttt{int a[5], *p;}，当执行了 \texttt{p=\&a[3];} 语句时，是将指针变量 \texttt{p} 指向了 \texttt{a} 数组的第\_\_\_\_\_\_个元素的地址。
\begin{enumerate}[label=\Alph*.]
    \item 2
    \item 3
    \item 4
    \item 5
\end{enumerate}

\newpage

%========================
\section{笔记 6：二维数组 + 指针进阶题}
\begin{coverbox}
\textbf{封面标题建议：}\\
专升本C语言｜二维数组配指针，很多人卡在这\\
\textbf{副标题建议：}\\
别死记，先把“谁指向谁”看清楚
\end{coverbox}

\subsection*{题目 1}
若有以下定义和语句：
\[
\texttt{int c[4][5], (*p)[5];}
\]
\[
\texttt{p=c;}
\]
能够正确引用数组元素 \texttt{c[1][2]} 的是
\begin{enumerate}[label=\Alph*.]
    \item \texttt{*(*(p+1)+2)}
    \item \texttt{(*p)[2]}
    \item \texttt{p+1[2]}
    \item \texttt{*(p+5)}
\end{enumerate}

\subsection*{题目 2}
设有说明 \texttt{int (*ptr)[10]}，其中的标识符 \texttt{ptr} 是
\begin{enumerate}[label=\Alph*.]
    \item 10 个指向整型变量的指针
    \item 指向 10 个整型变量的函数指针
    \item 一个指向具有 10 个整型元素的一维数组的指针
    \item 具有 10 个指针元素的一维指针数组，每个元素都只能指向整型量
\end{enumerate}

\subsection*{题目 3}
若有以下说明和语句：
\[
\texttt{int c[4][5], (*p)[5];}
\]
\[
\texttt{p=c;}
\]
则以下叙述中正确的是
\begin{enumerate}[label=\Alph*.]
    \item \texttt{p} 是指向整型变量的指针
    \item \texttt{p} 是指向整型变量的数组指针
    \item \texttt{p} 是指向函数的指针
    \item \texttt{p} 是数组名
\end{enumerate}

\newpage

%========================
\section{笔记 7：字符串和字符数组高频题}
\begin{coverbox}
\textbf{封面标题建议：}\\
专升本C语言｜字符串长度题，真的别再凭感觉做了\\
\textbf{副标题建议：}\\
这一组最容易在 \texttt{\textbackslash 0} 上翻车
\end{coverbox}

\subsection*{题目 1}
若有定义
\begin{lstlisting}
char s[10]="hello\0\t\\ ";
\end{lstlisting}
则 \texttt{printf("\%d",strlen(s));} 的输出结果是
\begin{enumerate}[label=\Alph*.]
    \item 10
    \item 9
    \item 7
    \item 6
\end{enumerate}

\subsection*{题目 2}
在下列给出的字符数组 \texttt{c} 中，它在内存中所占用的字节数是
\begin{lstlisting}
char c[]={"C language"};
\end{lstlisting}

\begin{enumerate}[label=\Alph*.]
    \item 10
    \item 11
    \item 12
    \item 13
\end{enumerate}

\subsection*{题目 3}
若有定义 \texttt{char* s="\\Name\\Address\textbackslash n"}，则指针 \texttt{s} 所指字符串长度为
\begin{enumerate}[label=\Alph*.]
    \item 19
    \item 15
    \item 18
    \item 说明不合法
\end{enumerate}

\newpage

%========================
\section{笔记 8：函数 + 传参与值传递}
\begin{coverbox}
\textbf{封面标题建议：}\\
专升本C语言｜函数参数到底改没改？3题看透\\
\textbf{副标题建议：}\\
会做这组题，函数传参就不容易混
\end{coverbox}

\subsection*{题目 1}
设有如下函数：
\[
\texttt{exam(float x)\{ printf("\%f",x*x); \}}
\]
则函数的类型为
\begin{enumerate}[label=\Alph*.]
    \item 与参数 \texttt{x} 的类型相同
    \item 是 \texttt{void}
    \item 是 \texttt{int}
    \item 无法确定
\end{enumerate}

\subsection*{题目 2}
以下程序的输出结果是
\begin{lstlisting}
#include <stdio.h>
void fun(char *c, int d)
{
    *c=*c+1; d=d+1;
    printf("%c, %c, ", *c, d);
}
void main()
{
    char a='A', b='a';
    fun(&b, a);
    printf("%c, %c\n", a, b);
}
\end{lstlisting}

\begin{enumerate}[label=\Alph*.]
    \item B, b, A, a
    \item a, B, B, b
    \item A, b, B, a
    \item b, B, a, B
\end{enumerate}

\subsection*{题目 3}
以下关于宏的叙述中正确的是
\begin{enumerate}[label=\Alph*.]
    \item 宏名必须用大写字母表示
    \item 宏定义必须位于源程序中所有语句之前
    \item 宏替换没有数据类型限制
    \item 宏调用比函数调用耗费时间
\end{enumerate}

\newpage

%========================
\section{笔记 9：结构体基础高频题}
\begin{coverbox}
\textbf{封面标题建议：}\\
专升本C语言｜结构体基础题，别只会背 \texttt{struct}\\
\textbf{副标题建议：}\\
这组题最适合做第一轮查漏补缺
\end{coverbox}

\subsection*{题目 1}
若有结构体定义：
\[
\texttt{struct Student \{ int num; char name[20]; float score; \};}
\]
则下列叙述正确的是
\begin{enumerate}[label=\Alph*.]
    \item \texttt{Student} 是结构体类型
    \item \texttt{num}、\texttt{name}、\texttt{score} 是结构体成员
    \item 结构体可以存放多种不同类型的数据
    \item 以上都正确
\end{enumerate}

\subsection*{题目 2}
下面关于结构体的说法中，错误的是
\begin{enumerate}[label=\Alph*.]
    \item 结构体可以包含不同类型的成员
    \item 结构体变量可以直接赋值
    \item 结构体变量可以比较大小
    \item 结构体可以嵌套定义
\end{enumerate}

\subsection*{题目 3}
结构体成员访问运算符是
\begin{enumerate}[label=\Alph*.]
    \item \texttt{.}
    \item \texttt{->}
    \item \texttt{*}
    \item \texttt{\&}
\end{enumerate}

\subsection*{题目 4}
若有定义：
\[
\texttt{struct Student \{ int num; char name[20]; \} stu, *p = \&stu;}
\]
则访问 \texttt{stu} 的 \texttt{num} 成员正确的是
\begin{enumerate}[label=\Alph*.]
    \item \texttt{stu.num}
    \item \texttt{p->num}
    \item \texttt{(*p).num}
    \item 以上都正确
\end{enumerate}

\newpage

%========================
\section{笔记 10：结构体 / 共用体 / 枚举辨析题}
\begin{coverbox}
\textbf{封面标题建议：}\\
专升本C语言｜struct、union、enum 总混？这组题专治混淆\\
\textbf{副标题建议：}\\
概念题最怕似懂非懂
\end{coverbox}

\subsection*{题目 1}
在 C 语言中，以下哪个选项表示结构体类型？
\begin{enumerate}[label=\Alph*.]
    \item \texttt{struct}
    \item \texttt{union}
    \item \texttt{enum}
    \item \texttt{class}
\end{enumerate}

\subsection*{题目 2}
在 C 语言中，以下哪个选项表示共用体类型？
\begin{enumerate}[label=\Alph*.]
    \item \texttt{struct}
    \item \texttt{union}
    \item \texttt{enum}
    \item \texttt{class}
\end{enumerate}

\subsection*{题目 3}
在 C 语言中，以下哪个选项表示枚举类型？
\begin{enumerate}[label=\Alph*.]
    \item \texttt{struct}
    \item \texttt{union}
    \item \texttt{enum}
    \item \texttt{class}
\end{enumerate}

\subsection*{题目 4}
下面关于共用体的说法中，错误的是
\begin{enumerate}[label=\Alph*.]
    \item 共用体所有成员共享同一段内存
    \item 共用体变量的大小是最大成员的大小
    \item 共用体可以同时存放多个成员的值
    \item 共用体变量初始化时只能给第一个成员赋值
\end{enumerate}

\end{document}